% ══════════════════════════════════════════════════════════════════════════════
\chapter{Moteur de recherche}
\label{chap:td6}
% ══════════════════════════════════════════════════════════════════════════════

Le sixième et dernier TD réalise l'intégration complète du système : le parser de requêtes
du TD5 est branché sur les index inversés du TD3, formant un moteur de recherche
opérationnel. La consigne demande également une évaluation empirique sur un jeu de test
tiré de l'annexe, et une interface utilisateur permettant d'interroger le corpus
interactivement. Ce TD illustre les notions de précision, rappel et pertinence au sens
du cours (chapitre~5, Modèle vectoriel et évaluation).

\section{Architecture et modèle de recherche}

Le moteur s'organise en quatre modules dans \texttt{src/adit\_corpus\_indexing/search/} :
\texttt{engine.py} orchestre le traitement complet d'une requête, \texttt{corpus\_reader.py}
maintient en mémoire les métadonnées de chaque article, \texttt{index\_loader.py} charge
les fichiers TSV, et \texttt{evaluator.py} mesure les performances. Le modèle implémenté
est un \textbf{modèle booléen classé} (\emph{ranked boolean}) : les mots-clés sont
pondérés par le score TF-IDF précalculé (formule~\ref{eq:tfidf}) issu des index TD3,
tandis que les filtres facettes opèrent en AND strict sur des ensembles d'identifiants.

\begin{lstlisting}[caption={Exécution d'une requête : scoring puis intersection des filtres},
  label={lst:engine_execute}]
def _execute(self, pq: ParsedQuery) -> list[SearchResult]:
    # 1. Score des mots-clés via index TF-IDF précalculé
    keyword_scores = self._search_keywords(pq.mots_cles, pq.operateurs, pq.zone)

    # 2. Ensembles de documents satisfaisant chaque filtre
    rubrique_docs = self._get_rubrique_docs(pq.rubrique) if pq.rubrique else None
    date_docs     = self._get_date_docs(pq.date_min, pq.date_max)
    image_docs    = self._corpus.ids_with_images() if pq.filtre_images else None

    # 3. Intersection AND de tous les filtres actifs
    result_ids = set(keyword_scores.keys()) or None
    for fset in (rubrique_docs, date_docs, image_docs):
        if fset is not None:
            result_ids = fset if result_ids is None else result_ids & fset
    ...
\end{lstlisting}

\section{Évaluation expérimentale}

Le ground truth a été constitué \textbf{indépendamment du moteur} par grep direct sur
les fichiers TSV et parsing XML, garantissant que les métriques reflètent des performances
réelles et non circulaires. Les mesures de précision $P$, rappel $R$ et F1 suivent les
définitions du cours :
\begin{equation}
  P = \frac{|TP|}{|TP|+|FP|}, \quad
  R = \frac{|TP|}{|TP|+|FN|}, \quad
  F_1 = \frac{2PR}{P+R}
  \label{eq:prf1}
\end{equation}

\begin{table}[ht]
\centering
\caption{Résultats d'évaluation sur 10 requêtes du ground truth (100 exécutions chacune)}
\label{tab:eval}
\begin{tabular}{clrrrrr}
\toprule
\textbf{\#} & \textbf{Requête (abrégée)} & \textbf{P} & \textbf{R} & \textbf{F1}
  & \textbf{Ret.} & \textbf{t (ms)} \\
\midrule
1  & Focus + images          & 1,000 & 1,000 & 1,000 & 56  & 1,04 \\
2  & Rubrique A lire         & 1,000 & 1,000 & 1,000 & 11  & 0,16 \\
3  & En direct des labos     & 1,000 & 1,000 & 1,000 & 41  & 0,36 \\
4  & Sans image              & 1,000 & 1,000 & 1,000 & 213 & 1,50 \\
5  & Titre $\ni$ nucléaire   & 1,000 & 1,000 & 1,000 & 4   & 0,63 \\
6  & Titre $\ni$ europe      & 1,000 & 1,000 & 1,000 & 1   & 0,13 \\
7  & Nov.\ 2011 + recherche  & 1,000 & 1,000 & 1,000 & 1   & 1,92 \\
8  & Focus + 08–09/2011      & 1,000 & 1,000 & 1,000 & 5   & 0,21 \\
9  & HE + systèmes embarqués & 1,000 & 1,000 & 1,000 & 1   & 0,12 \\
10 & Focus + 2014 + santé    & 1,000 & 1,000 & 1,000 & 3   & 0,26 \\
\midrule
\multicolumn{2}{l}{\textbf{Macro-moyenne}}
  & \textbf{1,000} & \textbf{1,000} & \textbf{1,000} & & \textbf{0,63} \\
\bottomrule
\end{tabular}
\end{table}

Le tableau~\ref{tab:eval} montre P\,=\,R\,=\,F1\,=\,1,000 sur toutes les requêtes,
avec un temps de réponse moyen de 0,63\,ms — soit 240 fois en dessous du seuil de
perception humaine. Ces métriques parfaites s'expliquent par la nature déterministe du
modèle booléen sur un corpus fermé : les filtres d'ensemble (rubrique, date, images)
sont exacts, et les requêtes mot-clé opèrent sur des racines Snowball sans ambiguïté.
Deux corrections ont été identifiées \emph{grâce} à l'indépendance du ground truth :
la normalisation des rubriques (variantes non prévues dans l'index) et le verbe de cadrage
\og{}voir\fg{} traité comme mot-clé actif, réduisant le rappel de la requête~Q8 de 1,0
à 0,4 avant correction.

\section{Interface utilisateur}

L'interface est une application web \textbf{Streamlit} dépassant la consigne minimale :
tri par pertinence, date croissante ou décroissante ; extrait contextuel (\emph{snippet})
de~200 caractères autour du premier mot-clé ; lien vers l'article HTML original ; panneau
d'\og{}Analyse de requête\fg{} décomposant la \texttt{ParsedQuery} ; panneau d'évaluation
interactif affichant P/R/F1 par requête avec explication pédagogique des métriques.
Le déploiement est reproductible via un \texttt{Dockerfile} inclus dans le dépôt.
